\begin{lstlisting}
int query(vector<int> tree[], int treeNode, int
start,
int end, int left, int right)
{
  if (start > right || end < left) {
    return 0;
  }
  if (start >= left && end <= right) {
    return tree[treeNode].end()
    - upper_bound(tree[treeNode].begin(),
    tree[treeNode].end(), right);
  }
  int mid = (start + end) / 2;
  int op1 = query(tree, 2 * treeNode, start, mid,
  left,
  right);
  int op2 = query(tree, 2 * treeNode + 1, mid + 1,
  end,
  left, right);
  return op1 + op2;
}
int main()
{
  int n; cin >> n;
  int arr[n];
  for(int i=0;i<n;i++) cin >> arr[i];
  int next_right[n];
  vector<int> tree[4 * n];
  unordered_map<int, int> ump;
  for (int i = n - 1; i >= 0; i--) {
    if (ump[arr[i]] == 0) {
      next_right[i] = n;
      ump[arr[i]] = i;
    }
    else {
      next_right[i] = ump[arr[i]];
      ump[arr[i]] = i;
    }
  }
  build(tree, next_right, 0, n - 1, 1);
  int q; cin >> q;
  while(q--){
    int l, r; cin >> l >> r;
    l-- , r--;
    cout << query(tree,1,0,n-1,l,r) << "\n";
  }
  return 0;
}
\end{lstlisting}

\textbf{Kth smallest value in a range (SegTree)}
\begin{lstlisting}
#define N (int)1e5
vector<int> seg[N];
void build(int ci, int st, int end,
pair<int, int>* B)
{
  if (st == end) {
    seg[ci].push_back(B[st].second);
    return;
  }
  int mid = (st + end) / 2;
  build(2 * ci + 1, st, mid, B);
  build(2 * ci + 2, mid + 1, end, B);
  merge(seg[2 * ci + 1].begin(), seg[2 * ci +
  1].end(),
  seg[2 * ci + 2].begin(), seg[2 * ci + 2].end(),
  back_inserter(seg[ci]));
}
\end{lstlisting}
